\documentclass[11pt]{article}
\usepackage[margin=1in]{geometry}
\usepackage{amsmath,amssymb,amsthm,mathtools}
\usepackage[T1]{fontenc}
\usepackage{lmodern,microtype}
\usepackage[colorlinks=true,linkcolor=blue,citecolor=blue,urlcolor=blue]{hyperref}
\setlength{\emergencystretch}{3em}
\newtheorem{theorem}{Theorem}
\newtheorem{lemma}[theorem]{Lemma}
\newtheorem{proposition}[theorem]{Proposition}
\theoremstyle{definition}
\newtheorem{definition}[theorem]{Definition}
\newtheorem{remark}[theorem]{Remark}
\newcommand{\ket}[1]{\lvert #1\rangle}
\newcommand{\bra}[1]{\langle #1\rvert}
\newcommand{\norm}[1]{\lVert #1\rVert}
\newcommand{\F}{\mathcal F}
\newcommand{\Id}{\mathrm I}
\newcommand{\poly}{\operatorname{poly}}
\newcommand{\SYT}{\operatorname{SYT}}
\newcommand{\cont}{\operatorname{cont}}
\DeclareMathOperator{\atan}{arctan}
\title{A direct-induction proof of a coherent\texorpdfstring{\\}{ }$\widetilde O(n^{5/2})$ symmetric-group QFT}
\author{Fourier decomposition project: research theorem note}
\date{23 October 2026}
\begin{document}
\maketitle

\begin{abstract}
We give a circuit-existence proof for the symmetric-group quantum Fourier
transform using $n^{5/2}\poly(\log(n+1),\log(1/\epsilon))$ elementary gates
and $n\poly(\log(n+1),\log(1/\epsilon))$ live qubits. The construction
prepares a column of the standard induction transform in near-linear cost,
erases its transversal label with amplitude $M^{-1/2}$, and amplifies the
entire induction unitary in $O(\sqrt M)$ rounds. Elementary bounded-product
arithmetic supplies the branching coefficients. Fixed sorting networks
supply tableau caches, complete promise reflections, and reversible
boundary updates. Coherent arithmetic, finite-gate synthesis, error
accumulation, and workspace cleanup are included. The proof does not use a
fast decomposition of the local A-matrix. This is an internally audited
research draft, not an externally reviewed result, a priority claim, or a
compiled fault-tolerant implementation.
\end{abstract}

\section{Statement and computational conventions}

A partition $\lambda\vdash m$ is a nonincreasing list of positive integers
with sum $m$. Its boxes have coordinates $(r,c)$, where $1\le c\le\lambda_r$,
and content $c-r$. A standard Young tableau fills the boxes by $1,\ldots,m$
in increasing order along every row and column. Write $d_\lambda=|\SYT(\lambda)|$.
We encode a tableau by its row word $P=(p_1,\ldots,p_m)$: label $j$ is in
row $p_j$. Its column is the within-row rank
\[
 c_j=|\{h\le j:p_h=p_j\}|.
\]
The endpoint shape is stored as a zero-padded list of row lengths, not in a
succinct or random-access data structure. Every word has $O(\log(n+1))$ bits.

Let $R_\lambda$ denote Young's real orthogonal representation, with positive
off-diagonal entries for adjacent transpositions. With the rightmost
permutation acting first, our Fourier convention is
\begin{equation}\label{eq:qft}
 \F_n\ket g=\sum_{\lambda\vdash n}\sum_{Q,R\in\SYT(\lambda)}
 \sqrt{\frac{d_\lambda}{n!}}\,
 \bra Q R_\lambda(g)\ket R\,\ket{\lambda,Q,R}.
\end{equation}
The shape register is redundant but is retained. Inputs use the unique
transversal word
\[
 g=t_{i_n}^{(n)}\cdots t_{i_1}^{(1)},\qquad
 t_i^{(M)}=(i,i+1,\ldots,M),\quad 1\le i_M\le M.
\]
The valid input and output subspaces have dimension $n!$. A circuit has an
arbitrary unitary extension on other physical bitstrings.

We count gates from a fixed finite universal set, for example Clifford+$T$,
including reversible arithmetic and rotation synthesis. There is no QRAM,
precomputed shape oracle, free real-number arithmetic, or postselection.
The public parameters are $n$ and an integer requested precision. All loops,
capacities, and wire addresses are fixed by these parameters. A clean-error
bound means operator norm on the valid input subspace with all auxiliary
registers zero, comparing against the ideal output with auxiliaries zero.
It therefore includes leakage into auxiliary registers and holds with an
arbitrary entangled reference system. Approximate cleanup means this norm
bound, not an assertion that every synthesized circuit ends with exactly
zero auxiliary bits.

\begin{theorem}[Main bound]\label{thm:main}
For $n\ge1$ and $0<\epsilon\le1/2$ there is a uniform finite-gate circuit
with clean-input isometry $\widetilde V_n$ such that
\[
 \norm{\widetilde V_n-(\F_n\otimes\ket0_{\rm work})}\le\epsilon/2.
\]
Its elementary-gate count and maximum number of live qubits are respectively
\[
 n^{5/2}\poly(\log(n+1),\log(1/\epsilon)),\qquad
 n\poly(\log(n+1),\log(1/\epsilon)).
\]
After discarding work, the channel is within diamond distance $\epsilon$
of the ideal channel on the input promise. The same dominant bound holds
for a one-line permutation input after a reversible polynomial-logarithmic
overhead, quadratic-time encoding conversion.
\end{theorem}

The $\widetilde O$ notation below suppresses polynomial factors in
$\log(n+1)$ and $\log(1/\epsilon)$; it is not uniform in unbounded exact-real
precision. The result is an approximation theorem, not an exact finite-gate
realization. The proof optimizes the power of $n$, not the logarithmic
exponents or practical constants.

\section{The standard induction transform}

Fix $\lambda\vdash M-1$. The notation $\mu\succ\lambda$ means that $\mu$
is obtained by adding one box to $\lambda$. If $P\in\SYT(\lambda)$, write
$P\to\mu$ for the tableau obtained by placing $M$ in that box. Define
\begin{equation}\label{eq:induction}
 (F_\lambda)_{(\mu,Q),(i,P)}=
 \sqrt{\frac{d_\mu}{M d_\lambda}}\,
 \bra Q R_\mu(t_i^{(M)})\ket{P\to\mu}.
\end{equation}
The standard multiplicity-free induction rule says that this is a unitary
map from $\mathbb C^M\otimes\mathbb C^{d_\lambda}$ to
$\bigoplus_{\mu\succ\lambda}\mathbb C^{d_\mu}$. In particular,
$\sum_{\mu\succ\lambda}d_\mu=M d_\lambda$. These representation-theoretic
facts, and the resulting QFT recurrence, are established ingredients, not
claims of this note; see Definition 3.6, Lemma 3.7, and Theorem 3.8 of
\cite{bgo}. We prove the recurrence's circuit and error costs below.

For clarity, the induction identity can be understood by decomposing
$\operatorname{Ind}_{S_{M-1}}^{S_M}R_\lambda$. In the coset basis it has
labels $(i,P)$. Each successor representation occurs once, its restricted
$\lambda$ sector has basis $P\to\mu$, and the matrix coefficients of the
intertwiner are \eqref{eq:induction}. Orthogonality fixes the normalization.
This is the only non-elementary representation-theoretic input to the
induction map itself.

\subsection{A common physical register space}

At a fixed stage $M$, the unchanged external register stores $\lambda$ in
$M$ padded row-length slots. There are a tag register with alphabet
$\{0,\ldots,M\}$, $M$ tableau row slots, an attenuation flag, and work
registers. The tag value $0$ and row value $0$ are reserved as blanks.
Define isometries $I_\lambda,J_\lambda$ into this same physical space:
\begin{align*}
 I_\lambda\ket{i,P}&=\ket{\lambda;i,P,0;0;0_{W}},\quad 1\le i\le M,\\
 J_\lambda\ket{\mu,Q}&=\ket{\lambda;0,Q;0;0_{W}},\quad\mu\succ\lambda.
\end{align*}
Here $0_W$ includes \emph{all} work used by the column unitary, including
coefficient, angle, column-cache, and sorting registers. The external shape
is part of the control, not an erased output. Define $I,J,F$ by direct sum
over valid $\lambda$. Their ranks vary by shape, but their physical wire
counts do not. Right tableaux and unused future input tags are spectators;
they are not tested for zero. Inactive tails of larger-capacity stage
registers are also untouched spectators, initially zero.

Two reflections will be required on these \emph{full} promises:
\begin{equation}\label{eq:reflect}
 R_I=2II^\dagger-\Id,\qquad R_J=2JJ^\dagger-\Id.
\end{equation}
Testing only the tag is not sufficient. Dedicated predicate helper registers
are kept separate from $W$-work and return exactly to zero after each
reflection. All identities involving these reflections are on the invariant
subspace where those dedicated helpers start at zero.

\section{Near-linear reversible tableau bookkeeping}\label{sec:bookkeeping}

\begin{lemma}[Fixed-schedule word operations]\label{lem:words}
For row words and padded shapes with $M$ slots, the following functions have
clean reversible XOR implementations using $M\poly(\log(n+1))$ gates and
work bits: within-row column ranks, endpoint row histograms, standard-tableau
validity, partition validity, and the condition $\mu\succ\lambda$.
Consequently the reflections \eqref{eq:reflect}, including zero tests of
$W$-work of size $M\poly(\log(n+1),\log(1/\eta))$, have cost and helper
space $M\poly(\log(n+1),\log(1/\eta))$.
\end{lemma}

\begin{proof}
Use a fixed bitonic sorting network on a power-of-two padding of $O(M)$
records. Its $O(M\log^2(M+1))$ compare-exchanges compare $O(\log(n+1))$-bit
keys and retain their comparison bits. Sentinels and original-position
labels are included in the keys, so every comparison is deterministic,
including on equal or malformed data. Reversing swaps and comparisons
restores the records and zeros the comparison bits. A scan uses a fresh
output word at each step and retains its old words until reversal. Modular
fixed-width arithmetic gives a total Boolean function on every input;
valid inputs fit in the chosen capacities.

For columns, sort records $(p_j,j)$ by row and then label. A segmented
prefix count gives the rank of each label within its row. Sorting by
original label places the resulting ranks at public output addresses.
For the histogram, insert one sentinel for each row $r=1,\ldots,M$ after
that row's data records. Its segmented count is the row length; sort these
sentinels by their public row labels for extraction. Check separately that
all input rows lie in $[1,M]$. The row histogram must be nonincreasing.
To test tableau validity, sort label records by their computed column and
then row. Within each column, check that rows are consecutive starting at
one and that labels increase. Within-row labels already increase by the
definition of a row word. This proves necessity and sufficiency, including
absence of holes. Partition validity is a fixed scan of row comparisons,
nonnegativity, and the total sum. Cover validity checks that the padded
rowwise differences are zero except for one $+1$, with both partitions
valid and the correct total sizes.

Each algorithm is made into a clean function oracle by computing its
intermediate history, copying the requested outputs by XOR, and reversing
the history. Thus it works even when the target XOR register is nonzero.
Malformed input is not an excuse to drop information: its computed
predicate is copied, then the same computation is reversed. No lookup by a
data-dependent row address is used.

For $R_I$, test all of the following: $\lambda$ is a partition of $M-1$;
the first $M-1$ row symbols encode a standard tableau whose endpoint is
exactly $\lambda$; the final row symbol is zero; $1\le i\le M$; and the
attenuation flag and every $W$-work bit are zero. For $R_J$, test that
$\lambda\vdash M-1$, all $M$ row symbols encode a standard tableau $Q$
whose endpoint $\mu$ covers $\lambda$, and the tag, flag, and every
$W$-work bit are zero. Untested right paths and future digits are
tensor-identity spectator factors in both projectors.
Compute the conjunction in
separate clean helper space, apply phase $-1$ precisely when it is false,
and uncompute. All this is exact reversible Boolean logic. The helpers
are \emph{not} borrowed from possibly dirty $W$-work.
\end{proof}

One bookkeeping point is essential later: the column cache is built by a
\emph{clean} XOR oracle before applying representation matrices. We do not
leave sorting history attached to the old tableau. After the tableau
changes, the same XOR oracle evaluated on the new tableau erases the
updated cache.

\section{Elementary coherent branching coefficients}\label{sec:weights}

Let $x_0>y_0>x_1>\cdots>y_{k-2}>x_{k-1}$ be the integer contents of
addable and removable boxes of $\lambda$, respectively. The standard
dimension-ratio identity gives, for the successor belonging to $x_i$,
\begin{equation}\label{eq:residue}
 p_i=\frac{d_{\lambda+\Box_i}}{M d_\lambda}
 =\frac{\prod_{j=0}^{k-2}|x_i-y_j|}{\prod_{h\ne i}|x_i-x_h|}.
\end{equation}
See also Lemma 2.4 of \cite{bgo}. Normalization follows directly by taking
the coefficient of $1/z$ in the partial fractions of
$\prod_j(z-y_j)/\prod_i(z-x_i)$. A scan of padded row lengths identifies
corners; a fixed stable compaction by sorting puts active corners first.
There are $k$ addable corners and $k-1$ removable corners, with
\[
 \frac{k(k-1)}2\le M-1,\qquad L:=x_0-x_{k-1}\le2M.
\]
Indeed $k-1$ distinct positive row lengths require at least
$1+\cdots+(k-1)$ boxes. Thus all arrays may be padded to the public
least power-of-two capacity
$K=2^{\lceil\log_2(\lceil\sqrt{2M}\rceil+1)\rceil}$, so $K^2=O(M)$.

\begin{lemma}[Bounded products and precision]\label{lem:precision}
For $k>1$, each probability in \eqref{eq:residue} has the factorization
\begin{equation}\label{eq:bounded}
 p_i=\prod_{j<i}\frac{y_j-x_i}{x_j-x_i}
       \prod_{j\ge i}\frac{x_i-y_j}{x_i-x_{j+1}},
\end{equation}
with every factor strictly between zero and one, and $p_i\ge L^{-2}$.
For each $i$ start with $N_i^{(0)}=2^b$ and, for a factor $a/d$, update
\[
 N_i^{(t+1)}=\left\lfloor N_i^{(t)}a/d\right\rfloor.
\]
Writing $N_i=N_i^{(k-1)}$, $T=\sum_iN_i$, and $\widehat p_i=N_i/T$,
the following bounds hold whenever the stated precision makes $T>0$:
\begin{align}
 0\le p_i-2^{-b}N_i&<(k-1)2^{-b},\label{eq:round}\\
 \norm{p-\widehat p}_1&\le2k(k-1)2^{-b},\label{eq:l1}\\
 \left\|\sum_i(\sqrt{p_i}-\sqrt{\widehat p_i})\ket i\right\|
 &\le\sqrt{2k(k-1)2^{-b}}.\label{eq:root}
\end{align}
In particular $b=O(\log(M+1)+\log(1/\eta))$ suffices for a clean
coherent preparation of $\sum_i\sqrt{p_i}\ket i$ with error at most $\eta$
and $M\poly(\log(n+1),\log(1/\eta))$ elementary gates and work bits.
The case $k=1$ is an exact deterministic state.
\end{lemma}

\begin{proof}
Regrouping denominators gives \eqref{eq:bounded}. On the left of $i$, use
$y_j\ge x_{j+1}+1$, and telescope all but the last positive integer
numerator, giving a lower bound $1/(x_0-x_i)$ when this side is nonempty.
On the right, use $y_j\le x_j-1$ and obtain $1/(x_i-x_{k-1})$ in the
same way. An empty side contributes one. This proves $p_i\ge L^{-2}$.
Each rounded update contracts the previous error and adds less than $2^{-b}$,
proving \eqref{eq:round}. If
$\Delta=1-2^{-b}\sum_iN_i$, then
$0\le\Delta\le k(k-1)2^{-b}$. Renormalization changes the dyadic vector
by exactly $\Delta$ in $\ell^1$, proving \eqref{eq:l1}.
The inequality $(\sqrt a-\sqrt b)^2\le|a-b|$ proves \eqref{eq:root}.

For example, take the public precision
\[
 b\ge\max\left\{\left\lceil\log_2(2(K-1)(2M)^2)\right\rceil,
             2\left\lceil\log_2(4K/\eta)\right\rceil\right\}.
\]
The weight loss is at most $(K-1)2^{-b}\le1/[2(2M)^2]\le p_i/2$,
so every active $N_i$ is positive. Weight-rounding state error is less
than $\eta/2$. Inactive leaves are exactly zero. Throughout an update
$N_i\le2^b$ and its product with a numerator is at most $2^bL$; there
are no exponentially large intermediate products. Fixed multiplication
and division use polynomially many gates in $b+\log(M+1)$. Compute one
product history, copy its final word, and reverse that history before
moving to the next $i$. This uses $O(K^2)$ updates and $O(K)$ words of
history at a time, up to polynomial word-arithmetic overhead.

Build the binary tree of the integer masses $N_i$. Each mass is at most
$K2^b$. At a node with child masses $A,B$, use the real rotation of angle
$\theta=\operatorname{atan2}(\sqrt B,\sqrt A)$, controlled by the already
prepared prefix. Zero children use exact identity or bit-flip endpoints;
a zero parent uses identity. There are $K-1$ fixed nodes. Compute each
angle to error at most $\eta/(8K)$ and synthesize its controlled rotation
to operator error at most $\eta/(8K)$. Telescoping the gates gives error
at most $\eta/4$, leaving slack in the claimed bound. Each node's prefix
condition uses at most $\log_2 K$ bits: compute its equality flag in clean
helper bits, apply the controlled rotation, and uncompute the flag. This
adds only polynomial-logarithmic cost per node, with reusable helper space.

These angle computations need no oracle: evaluate
$z=\sqrt{\min(A,B)/\max(A,B)}\in[0,1]$ by integer division and a fixed
integer-square-root routine to enough dyadic places. For $z\le1/2$ the
alternating series for $\atan z$ converges geometrically. For $z>1/2$ use
$\atan z=\pi/4-\atan((1-z)/(1+z))$; the new argument is at most $1/3$.
Machin's formula computes $\pi/4=4\atan(1/5)-\atan(1/239)$ to the same
accuracy. Put $c=O(\log(K/\eta))$. Evaluate the $O(c)$ terms by fixed-point
Horner evaluation, rounding each operation at $c+O(\log c)$ fractional
bits; these guard bits bound the accumulated arithmetic error by $O(2^{-c})$.
Exact rational powers are useful reference certificates but are not the
claimed bounded-word circuit implementation.
The square-root input can be evaluated to twice the desired output bits,
using $|\sqrt u-\sqrt v|\le\sqrt{|u-v|}$. If $B>A$, subtract the small
angle from $\pi/2$. Every step has polynomial bit complexity.

The $O(c)$ angle bits control $O(c)$ fixed dyadic rotations. Assign each
bit-controlled rotation synthesis error $O(\eta/(Kc))$; their sum is
$O(\eta/K)$ per tree node. Finite-gate synthesis has
polynomial-logarithmic word length \cite{sk}. For classical compilation
in a finite-precision bit model we use Theorem 4.5 of \cite{mw}, rather
than interpreting an arithmetic-operation count as a bit-complexity bound.
The approximation here is of the actual controlled gate in matrix norm,
not merely a one-qubit operation up to an uncontrolled global phase.
For instance a controlled real rotation decomposes into two CNOTs and a
half-angle rotation paired with its exact circuit adjoint; its irrelevant
single-qubit synthesis phase then cancels. Prefix and angle-bit conjunctions
are computed and uncomputed by exact Boolean circuits. Their cost is
included in the polynomial-logarithmic factor.
Compute, rotate, and uncompute the angle
scratch. Tree masses, weight words, and contents are erased from the
unchanged shape after preparation. The loop bounds use $K$, not the
data-dependent $k$. Dummy factors are one, dummy leaves zero, and all
invalid denominators are replaced by one in totalized arithmetic. A
shape-validity guard may select a fixed default state on invalid shapes.
All controlled maps therefore have a unitary extension, with no need to
specialize or enumerate circuits over partitions.
\end{proof}

To obtain the actual added row rather than its corner ordinal, use a fixed
scan of the $K$ corner records: controlled on ordinal equality, XOR the
selected row into the final blank tableau slot. A second fixed scan,
controlled on equality with that unique row, erases the ordinal. This is
reversible on the prepared support. On the complement it is still a
composition of reversible XOR gates. The parameter table is then erased
using the unchanged $\lambda$. Thus Lemma~\ref{lem:precision} implements
the branching column map
\begin{equation}\label{eq:B}
 B:\ket{P,0}\longmapsto
 \sum_{\mu\succ\lambda}\sqrt{d_\mu/(M d_\lambda)}\ket{P\to\mu}
\end{equation}
with the same clean-error and cost bounds. A fixed-angle approximate
circuit can leak into padded states; its entire leakage is included in the
isometry error, not assumed to vanish exactly.

\section{The selected transversal in near-linear cost}\label{sec:generators}

Write $s_j=(j,j+1)$ and $d=\cont(j+1)-\cont(j)$ in a valid tableau.
Young's orthogonal rule is
\[
 R_\mu(s_j)\ket P=d^{-1}\ket P+
 \sqrt{1-d^{-2}}\ket{s_jP}.
\]
When $|d|=1$ the second term is absent and the action is the exact sign $d$.
Otherwise $s_jP$ is standard and its axial distance is $-d$.

\begin{lemma}[Cached generators]\label{lem:generators}
With a correct column cache, one public-index Young generator, also
controlled on a public comparison with the tag, can be implemented to
operator error $\zeta$ on the valid-tableau/cache subspace in
$\poly(\log(n+1),\log(1/\zeta))$ gates and work bits. The ideal operation
preserves that subspace and updates the column cache correctly.
\end{lemma}

\begin{proof}
Only the two records $(p_j,c_j)$ and $(p_{j+1},c_{j+1})$ are accessed.
For $|d|>1$, compute $f=[d<0]$ and swap both records if $f=1$, thereby
canonicalizing the pair to a positive distance $D=|d|$. Retain $f$; this
step must not discard which member of the pair was present. On $f$ apply
\[
 A_D=\begin{pmatrix}1/D&\sqrt{1-D^{-2}}\\
             \sqrt{1-D^{-2}}&-1/D\end{pmatrix}
       =R_y(2\arccos(1/D))Z.
\]
Compute and erase $D$ and the angle while the canonical records are
unchanged. Undo the record swap controlled on the now-superposed $f$.
Finally XOR $[d_{\rm final}<0]$ into $f$, erasing it on both branches.
The $|d|=1$ branch applies the exact sign; its guard is invariant and can
be erased. On invalid inputs use safe total arithmetic and a reversible
extension of this gate sequence, without making a promise of cleanup.

The swapped column cache is correct because $s_jP$ interchanges the boxes
containing $j,j+1$; all other box coordinates are unchanged. If
$\theta=\arccos(1/D)$, then
$\norm{A_D(\theta)-A_D(\widetilde\theta)}
=2|\sin((\theta-\widetilde\theta)/2)|\le|\theta-\widetilde\theta|$.
Compute $\theta$ by the same bounded-ratio square-root/arctangent routine
as above, or via $\pi/2-\atan(1/\sqrt{D^2-1})$. The latter ratio is at
most $1/\sqrt3$; a further range reduction gives a geometrically
convergent series. Word precision $O(\log(n+1)+\log(1/\zeta))$ suffices.
Arithmetic is uncomputed and controlled dyadic rotations are synthesized
with the claimed polynomial-logarithmic cost. For a $c$-bit angle word,
assign each of its $O(c)$ controlled constant rotations error $O(\zeta/c)$,
using the same phase-aware construction as above. The axial-distance
sectors are orthogonal controls, so the error is the maximum sector error,
not a sum over tableaux. Only the subsequent $M-1$ generator applications
are summed; $\zeta=O(\eta/M)$ suffices for the selected transversal.
\end{proof}

Build the clean column cache of the complete $M$-box tableau, then, in
chronological order $j=M-1,M-2,\ldots,1$, apply this generator controlled
on $[i\le j]$. Since $t_i=s_i s_{i+1}\cdots s_{M-1}$, this order implements
\[
 V:\ket{i,Q}\longmapsto\ket i R_\mu(t_i)\ket Q.
\]
The tag zero case can act as identity. The cache is updated at each swap,
and is erased by the clean column XOR oracle evaluated on the final
tableau, not by reversing stale pre-generator history. With per-generator
error $O(\eta/M)$, the whole selected operation has clean error $O(\eta)$
and cost $M\poly(\log(n+1),\log(1/\eta))$. This error statement follows
by telescoping on ideal intermediate valid-cache subspaces; it does not
require approximate intermediate states to satisfy exact validity.

\section{Column erasure and full-unitary amplification}\label{sec:oaa}

Let $H_M$ be the specific tag unitary defined by a fixed binary preparation
tree, with $H_M\ket0=M^{-1/2}\sum_{i=1}^M\ket i$. The tree has $O(M)$
leaves; approximating each of its fixed rotations gives a whole-unitary
operator-norm approximation of $H_M$, and hence of its adjoint, to error
$\eta$ in $M\poly(\log(n+1),\log(1/\eta))$ gates.
No stronger cost bound for this elementary uniform
state is needed. Define the ideal unitary
\[
 W_0=(H_M^\dagger\otimes\Id)VB.
\]
Projecting onto tag zero and all work zero gives the exact block identity
\begin{equation}\label{eq:block}
 J^\dagger W_0 I=F/\sqrt M.
\end{equation}
Projection here specifies a matrix block; it is not a measurement.

Choose public values
\begin{equation}\label{eq:attenuate}
 r_M=\lceil\sqrt M\rceil,\qquad
 \theta_M=\frac{\pi}{4r_M+2},\qquad a_M=\sqrt M\sin\theta_M.
\end{equation}
Since $a_M\le\pi/4<1$, a one-qubit flag rotation with zero amplitude $a_M$
is legal. Apply it in addition to $W_0$ and call the result $W$.
The target tests flag zero, giving
\begin{equation}\label{eq:attenuated}
 J^\dagger W I=\sin\theta_M F.
\end{equation}
Computing and synthesizing this public angle has polynomial-logarithmic
cost. Special-case $M=1$ as a trivial induction if desired.

\begin{lemma}[Two-promise OAA and its error]\label{lem:oaa}
Suppose $I,J$ are isometries, $F$ is unitary between their equal-dimensional
domains, and a unitary $W$ obeys $J^\dagger WI=\sin\theta F$ with
$\theta=\pi/(4r+2)$ for an integer $r\ge1$. Then
\begin{equation}\label{eq:exactoaa}
 (-WR_IW^\dagger R_J)^r WI=JF.
\end{equation}
If $\widetilde W$ is unitary and
$\norm{(\widetilde W-W)I}\le\eta$, keeping the exact promise reflections
gives clean output error at most $(4r+1)\eta$.
\end{lemma}

\begin{proof}
Put $A=WI$, $G=JF$, $s=\sin\theta$, $c=\cos\theta$, and
$C=(A-sG)/c$. The block identity implies $G^\dagger C=0$ and
$C^\dagger C=\Id$, with $A=sG+cC$. In particular $J^\dagger C=0$.
On each corresponding two-dimensional plane,
$-WR_IW^\dagger R_J=-(2AA^\dagger-\Id)(2JJ^\dagger-\Id)$ sends
\[
 G\mapsto\cos(2\theta)G-\sin(2\theta)C,\qquad
 C\mapsto\sin(2\theta)G+\cos(2\theta)C.
\]
Its $r$th power applied to $A$ is
$\sin((2r+1)\theta)G+\cos((2r+1)\theta)C=G$.
The minus sign in \eqref{eq:exactoaa} is included; it avoids an unwanted
phase convention. This is the usual oblivious-amplification geometry
\cite{oaa}, written for distinct complete source and target promises.

For the error bound, set $\widetilde A=\widetilde W I$. Then
$\norm{\widetilde A\widetilde A^\dagger-AA^\dagger}\le2\eta$,
so the two image reflections differ by at most $4\eta$. Telescoping
the $r$ rounds and the initial isometry gives $(4r+1)\eta$.
Only the action of $W$ on the whole source promise must be approximated;
no bound on its arbitrary unitary completion is needed. The implemented
inverse is the exact adjoint circuit $\widetilde W^\dagger$, not a
separately synthesized inverse. If the two reflections themselves had
global errors $\rho_I,\rho_J$, one would add $r(\rho_I+\rho_J)$.
Here the Boolean promise reflections are exact.
\end{proof}

Lemmas~\ref{lem:words}--\ref{lem:generators}, the tag preparation, and the
flag rotation together construct $\widetilde W$ with source-isometry
error $\eta$ at cost $M\poly(\log(n+1),\log(1/\eta))$. Split this budget
among branching preparation, the selected generators, tag preparation,
and the flag. At fixed $\lambda$ the source tableau is a spectator during
branching preparation, so its state-error bound is an operator-isometry
bound. Direct sum over $\lambda$ takes a maximum of these errors, not a
sum over the number of partitions. The public capacities and number of
amplification rounds are the same for every coherent $\lambda$ sector.

There are $2r_M+1$ calls to $\widetilde W$ or its adjoint and $2r_M$
promise reflections. Thus a complete clean induction step costs
\begin{equation}\label{eq:stagecost}
 M^{3/2}\poly(\log(n+1),\log(1/\eta))
\end{equation}
gates and $M\poly(\log(n+1),\log(1/\eta))$ work bits. No branch is
postselected and no shape-dependent angle table is supplied for free.

\section{Outer boundary, Fourier recurrence, and total error}

At stage $M$ the previous Fourier registers have shape $\lambda$ and
tableaux $P,R\in\SYT(\lambda)$. Apply the amplified $F_\lambda$ to
$(i_M,P)$, leaving $R$ unchanged. On the ideal output, compute the new
shape $\mu$ from $Q$, and compute the unique row $a$ where $\mu-\lambda$
equals one. Append $a$ to $R$ to form $R\to\mu$.

This boundary map is clean and reversible in $\widetilde O(M)$ cost.
For example, keep the newly computed $\mu$, transfer $a$ to the blank
last right-path slot, and erase temporary difference work while $\lambda$
is still available. Then XOR the reconstructible old shape
$\mu-\Box_a$ into the old $\lambda$ register. It becomes zero. Swap the
new and old shape registers, retaining only $\mu$. The former new-shape
register is now zero. All row scans, histograms, and compute-copy-reverse
operations have fixed schedules. On malformed inputs these same XOR
operations still define a reversible circuit; residuals need not vanish.
On approximate inputs their norm is charged to the preceding error.

To verify \eqref{eq:qft} inductively, write $g=t_i^{(M)}h$ with
$h\in S_{M-1}$. The old coefficient times the induction coefficient is
\[
 \sqrt{\frac{d_\lambda}{(M-1)!}}\,
 \bra P R_\lambda(h)\ket R\,
 \sqrt{\frac{d_\mu}{M d_\lambda}}\,
 \bra Q R_\mu(t_i)\ket{P\to\mu}.
\]
Sum over $P$ and use the adapted restriction of $R_\mu$ to the
$\lambda$ sector. This is
$\sqrt{d_\mu/M!}\bra Q R_\mu(t_i h)\ket{R\to\mu}$,
exactly the next Fourier coefficient. Starting from the empty tableau
proves the whole recurrence, including the right-path update.

\begin{proof}[Proof of Theorem~\ref{thm:main}]
Use the public per-column error budget
\begin{equation}\label{eq:budget}
 \eta=\frac{\epsilon}{2n(4\lceil\sqrt n\rceil+1)}.
\end{equation}
By Lemma~\ref{lem:oaa}, the sum of clean induction errors is at most
$n(4\lceil\sqrt n\rceil+1)\eta=\epsilon/2$. Exact boundary maps do not
increase it. This hybrid argument evaluates each individual error on the
ideal valid input to that stage; actual accumulated leakage is carried
unitarily and requires no intermediate projection or measurement.
All word precisions are $O(\log(n+1)+\log(1/\epsilon))$, up to guard
bits of the same asymptotic size. Summing \eqref{eq:stagecost} gives
\[
 \sum_{M=1}^n O(M^{3/2})\poly(\log(n+1),\log(1/\epsilon))
 =n^{5/2}\poly(\log(n+1),\log(1/\epsilon)).
\]
Scratch space is reused across calls and stages, and all persistent
transversal, shape, and tableau words occupy $O(n\log(n+1))$ bits.
This proves the live-space bound. Reversible multiplication, division,
sorting, and finite-precision series have fixed polynomial-time Boolean
circuits. Fixed dyadic rotations are synthesized uniformly in polynomial
time in their precision in the bit model \cite{mw}. Hence the circuit
family is uniform. Explicit conservative logarithmic exponents and the
classical compilation bound are given in Section~\ref{sec:explicitcosts}.
Finally, clean isometry error at most $\epsilon/2$ implies channel
diamond error at most $\epsilon$, also after tracing out work.
\end{proof}

\subsection{Input-encoding cost}

A one-line permutation $g(1),\ldots,g(n)$ can be converted within the
same dominant bound without relying on a fast encoding theorem. At the
current size $M$, move the last image $i_M=g(M)$ into a blank digit slot,
leaving the last image slot zero. On each of the remaining images apply
$(t_{i_M}^{(M)})^{-1}$: the value $i_M$ maps to $M$, larger values decrease
by one, and smaller values are unchanged. This is a genuine permutation
of $[1,M]$, not a noninvertible conditional decrement. On valid inputs
the remaining images never equal $i_M$, so they now form a permutation
of $[1,M-1]$. Continue downwards. Reversal reconstructs the original
permutation from the saved digits, without retaining earlier arrays.
Public scans make $O(n^2)$ controlled word permutations, each implemented
by reversible comparisons and arithmetic in polynomial word cost and
space. For example, compute its piecewise bijection into a temporary word,
exchange the words, and erase the temporary by its inverse piecewise
bijection. Invalid digit values can select identity and out-of-range
image values can be fixed. Thus the full conversion is a reversible
unitary extension using $O(n\poly(\log(n+1)))$ space, and its
$O(n^2\poly(\log(n+1)))$ gate cost is lower-order than the QFT cost.

\section{Explicit precision and resource audit}\label{sec:explicitcosts}

The following refinement makes the suppressed costs checkable. Define
\[
 \Lambda=2+\lceil\log_2(n+1)\rceil+\lceil\log_2(1/\epsilon)\rceil,
\]
and use $\eta$ from \eqref{eq:budget}. These are sufficient bounds, not
optimized logarithmic exponents or predictions of practical circuit size.

\begin{proposition}[Explicit resource bounds]\label{prop:explicit}
The construction can be chosen to use
\[
 O(n^{5/2}\Lambda^5)\text{ elementary Clifford+$T$ gates},\qquad
 O(n\Lambda^3)\text{ live logical qubits}.
\]
A deterministic classical algorithm emits the circuit in
$n^{5/2}\poly(\Lambda)$ bit operations. Shape-dependent angle arithmetic
is performed coherently inside the circuit and is included in the gate
bound, not charged only to this public classical compilation.
\end{proposition}

\begin{proof}
At level $M\ge2$, take the public corner capacity $K$ above and set
\begin{align*}
 b&=\max\left\{\left\lceil\log_2(2(K-1)(2M)^2)\right\rceil,
           2\left\lceil\log_2(16K/\eta)\right\rceil\right\},\\
 c&=\left\lceil\log_2\frac{256(M+K+1)}\eta\right\rceil,\qquad
 q=c+\lceil\log_2(c+2)\rceil+12.
\end{align*}
Every parameter is $O(\Lambda)$. The positivity estimate of
Lemma~\ref{lem:precision} applies, and the weight-state error is at most
$\eta/8$.

Here is a fully bounded-word angle algorithm. For positive integer masses
$A,B$, find the largest integer $v$ with
\[
 v^2\max(A,B)\le 2^{2q}\min(A,B)
\]
by $q+1$ fixed bit trials. Thus $r=v2^{-q}$ approximates the desired
square root from below within $2^{-q}$. When $r>1/2$, set
$w=\operatorname{floor}_q((1-r)/(1+r))$ and later subtract
$\atan w$ from $\pi/4$; otherwise set $w=r$. In either case
$0\le w\le1/2$. Evaluate
\[
 \atan w\approx w\sum_{j=0}^{q+1}\frac{(-1)^j(w^2)^j}{2j+1}
\]
by signed Horner evaluation, rounding each multiplication and each public
coefficient to $q$ fractional bits. Two additional integer/sign bits
suffice for the accumulators. Compute the public $\pi/4$ by the same
routine and Machin's formula, with the rational arguments rounded down;
double it for $\pi/2$. Complement the angle when $B>A$.

For completeness, put $u=2^{-q}$. The rounded Horner multiplier is at
most $1/4$, the exact intermediate polynomials have magnitude at most
$4/3$, and each coefficient, product, and multiplier rounding contributes
at most a constant times $u$. The resulting error recurrence is bounded
by $e_j\le e_{j+1}/4+4u$. The series remainder is below $u$.
Root rounding contributes at most $u$ because $|\atan'|\le1$; range
reduction adds at most $u$, using the exact identity at the rounded $r$.
Machin's coefficients have absolute sum five; complementing adds at most
two copies of the public constant error. Summing these estimates is
bounded conservatively by
\[
 64(q+1)2^{-q}<2^{-c}.
\]
Zero parents select identity, and one-sided masses select exact endpoint
gates. No ratio lower bound or exponentially large exact rational powers
are needed. Public tag angles use the same routine. For attenuation,
an additional $O(\log M)$ guard allowance compensates for multiplication
by $\sqrt M$; $a_M\le\pi/4$ keeps the final inverse-cosine derivative
bounded. Elementary series on bounded arguments compute these public
constants in polynomial bit time to angle error $2^{-c}$.

Let $D=2^{\lceil\log_2(M+1)\rceil}$. There are at most
$N=K+M+D-2$ angle calls: $K-1$ branching nodes, $M-1$ generators,
$D-1$ tag nodes, and one attenuation. Thus their summed arithmetic error
is below $\eta/64$. There are at most
\[
 S=(K+M-2)(q+2)+D
\]
public rotation slots after expanding each computed angle into its bits
and allowing for endpoints. Assign each slot synthesis error
$\delta=\eta/(16S)$. Its total contribution is at most $\eta/16$.
Together with weight rounding this is strictly less than $\eta$, with
slack for the public-constant evaluation. Controlled rotations use the
phase-cancelling exact-adjoint construction already described, synthesizing
each half rotation to error at most $\delta/2$. An
uncontrolled rotation can use that construction with one clean auxiliary
temporarily set to one and restored. Overall minus signs can be realized
exactly by the Clifford identity $XZXZ=-\Id$ at constant cost.

On $O(\Lambda)$-bit words, addition and comparisons cost $O(\Lambda)$
Boolean gates, and schoolbook multiplication and restoring division cost
$O(\Lambda^2)$. The $O(\Lambda)$ root trials and Horner steps therefore
cost $O(\Lambda^3)$ per angle. Retaining a Boolean computation history
uses at most a constant number of fresh bits per gate; compute--copy--reverse
only changes constants. Sorting $O(M)$ records requires
$O(M\log^2 M)$ comparisons of $O(\log n)$-bit records, hence
$O(M\Lambda^3)$ gates and retained history. All $K(K-1)=O(M)$ weight
updates cost $O(M\Lambda^2)$. Even retaining every weight-update history
simultaneously would stay within this bound.

Finite synthesis can use $O(\log^4(1/\delta))=O(\Lambda^4)$ gates per
public rotation. For a bit-model uniform version, use \cite{mw}; its
single-qubit construction has five emitted recursive pieces per precision
improvement $\varepsilon_{\ell}=O(\varepsilon_{\ell-1}^{3/2})$, giving
word-length exponent $\log_{3/2}5<4$, and polynomial bit compilation
time. Phase-normalized single-qubit generators may be replaced by their
physical Clifford+$T$ representatives inside the cancelling sandwich.
The $O(\Lambda)$ rotation bits per angle therefore give $O(\Lambda^5)$
gates per angle. Synthesis words act on constant-size registers and do
not require a new history qubit per emitted quantum gate.

It follows that one column uses $O(M\Lambda^5)$ gates and
$S_W=O(M\Lambda^3)$ work bits. A complete promise reflection tests
$S_W$ work bits by a linear conjunction, in addition to the
$O(M\Lambda^3)$ tableau predicates. Its separate clean helpers also
occupy $O(M\Lambda^3)$ bits. The helpers are not part of $W$ and do
not recursively need their own validity test. Amplification repeats gates,
not physical registers; its $O(\sqrt M)$ repetitions give the claimed
sum over $M$. Exact outer boundaries cost $O(M\Lambda^3)$ per stage.
The hybrid proof permits reuse even when actual scratch has small leakage.

All public angles are polynomial-bit computable, the finite synthesis
algorithm has polynomial bit cost, and the fixed Boolean schedules can be
emitted in polynomial time per word operation. Including wire addresses
and writing the complete output gives $n^{5/2}\poly(\Lambda)$ classical
bit time. This statement does not claim that the original synthesis
arithmetic-operation exponent is its bit-complexity exponent.
\end{proof}

\section{Resource ledger and scope of the result}

With $b=\poly(\log(n+1),\log(1/\epsilon))$, the following deliberately
unoptimized ledger suppresses only polynomial factors in $b$.
\begin{center}
\begin{tabular}{lcc}
\hline
Operation at size $M$ & Elementary gates & Live work bits\\
\hline
Corners and $K^2=O(M)$ bounded products & $O(M\poly(b))$ & $O(M\poly(b))$\\
Branching tree and ordinal/row conversion & $O(\sqrt M\poly(b))$ & $O(M\poly(b))$\\
Cache construction and cleanup & $O(M\poly(b))$ & $O(M\poly(b))$\\
Selected $M-1$ Young generators & $O(M\poly(b))$ & $O(M\poly(b))$\\
Tag preparation and attenuation & $O(M\poly(b))$ & $O(M\poly(b))$\\
Each complete promise reflection & $O(M\poly(b))$ & $O(M\poly(b))$\\
Amplified induction, $O(\sqrt M)$ rounds & $O(M^{3/2}\poly(b))$ & $O(M\poly(b))$\\
Outer boundary update & $O(M\poly(b))$ & $O(M\poly(b))$\\
\hline
Complete $S_n$ QFT & $O(n^{5/2}\poly(b))$ & $O(n\poly(b))$\\
\hline
\end{tabular}
\end{center}

The coefficient-computation cost and the number of representation rotations
are both charged inside every column call, and those calls are charged
inside amplification. Classical compilation concerns only the public
fixed schedule and fixed dyadic angles. Shape-dependent arithmetic runs
reversibly in the quantum circuit; it is not an offline preprocessing
assumption. We make no claim that this implementation is practically
smaller than quadratic-coefficient alternatives at accessible sizes.

The published comparison point in \cite{bgo} is a $\widetilde O(n^3)$
upper bound in its stated approximation convention. The present draft's
formal exponent is smaller by $1/2$, a factor $\sqrt n$ before suppressed
logarithms. This comparison is not a priority determination. The present
argument does not prove that $5/2$ is optimal among all QFT circuits.
It also does not improve the exact Givens count of a local A-matrix.
It bypasses that decomposition in the global induction algorithm.

\subsection{Evidence, reproducibility, and remaining work}

The proof is analytic; finite tests do not establish its asymptotic claims.
The repository has independent dense tests of induction matrices, their
two-promise OAA completions, and the full small symmetric-group transform
in \texttt{alternative\_qft\_probe.py}. The new
\texttt{direct\_induction\_precision.py} supplies exact-rational rounding
and interval-angle checks. The new
\texttt{direct\_induction\_bookkeeping.py} models reversible fixed-network
records, predicates, and cleanup. The tests in
\texttt{test\_direct\_induction\_contract.py} check the cached-generator
canonicalization, sign and cache erasure, selected-transversal order,
promise padding, amplification perturbations, and the outer boundary.
These are reference and contract verifiers.

The separate \texttt{direct\_induction\_gate\_probe.py} now compiles a
complete $M=3$ induction step, including shape-dependent arithmetic,
complete promise reflections, and cleanup, into a lazily emitted
Clifford+$T$ program. All six valid source columns, across both shapes,
were compared with the independent ideal induction matrix. With nine
fractional weight and angle bits, the measured clean operator error is
$0.00505832$. The unoptimized program uses $17{,}138{,}782$ elementary
gates, including $7{,}778{,}300$ $T/T^\dagger$ gates, and $76{,}009$
logical qubits, excluding error correction. Two global-sign annotations
are excluded from the elementary-gate count. The saved report is
\texttt{data/direct\_induction\_gate\_M3\_b9\_c9.json}.
The simulation evaluates emitted reversible blocks and matrices multiplied
from actual finite gate words; its pruning bound does not certify
floating-point roundoff. This is a small implementation check, not a
compiled full QFT, a practical speedup, or evidence establishing asymptotics.

Remaining implementation work includes reducing the large arithmetic and
synthesis constants and assembling the full outer recurrence at small $n$.
Independent external review of the theorem and a separate systematic
priority search are also needed before a publication or novelty claim.
The note was developed with
AI-assisted derivation, coding, and internal adversarial checks in this
project; those checks are not independent peer review. No external human
review or publication acceptance is asserted.

\begin{thebibliography}{9}
\bibitem{bgo}
Carli Bruinsma, Dmitry Grinko, and Maris Ozols.
\emph{Quantum Fourier transform for the symmetric group}.
arXiv:2608.28569, version 2, 2026.
\url{https://arxiv.org/abs/2608.28569}.

\bibitem{oaa}
Dominic W. Berry, Andrew M. Childs, Richard Cleve, Robin Kothari, and
Rolando D. Somma.
\emph{Exponential improvement in precision for simulating sparse Hamiltonians}.
arXiv:1312.1414, 2014 version; Proceedings of STOC 2014.
\url{https://arxiv.org/abs/1312.1414}.

\bibitem{sk}
Christopher M. Dawson and Michael A. Nielsen.
\emph{The Solovay--Kitaev algorithm}.
Quantum Information \& Computation 6, 81--95 (2006).
\url{https://arxiv.org/abs/quant-ph/0505030}.

\bibitem{mw}
Dieter van Melkebeek and Thomas Watson.
\emph{Time-Space Efficient Simulations of Quantum Computations}.
Theory of Computing 8, 1--51 (2012), Theorem 4.5 and Sections 4.3.2--4.3.3.
\url{https://theoryofcomputing.org/articles/v008a001/}.
\end{thebibliography}
\end{document}
